\section{Proof of Lemmas}\label{sec:pf_lemmas}

Throughout this section, we use $b$ for the batch index and $B$ for the total number of batches.

\subsection{Proof of Lemma~\ref{lemma::throughput gap, single}}\label{appendix_lemma::proof of throughput gap, single}
Write \(n\) for the critical single-type threshold, so \(X^b\sim\mathrm{Poisson}(n)\). Taking expectations on both sides of the state transition equation gives
$$\ex{}{W^{b+1}}=\ex{}{W^b}+n - n \mathbb{P}(W^b+X^b\geq n).$$
Summing from \(b=0\) to \(B-1\), and noting that \(W^0=0\), we obtain
$$\ex{}{W^B}=n B - n \ex{}{B-B_{\text{stuck}}}=n\cdot \ex{}{B_{\text{stuck}}}.$$

\subsection{Proof of Lemma~\ref{lem:coupled_process, single}}\label{appendix_lemma::proof of coupled_process, single}
We prove by induction on the batch index \(b\). The coupled process \(\{\tilde{W}^b\}\) starts with a safety buffer of \(2n\), compared to the real process starting at 0, ensuring it remains above \(W^b+n\) throughout. For the inductive step, we consider two cases based on whether the system can process a batch or experiences a stuck iteration.

\paragraph{Base case.} For \(b=0\), \(\tilde{W}^0=2n\geq n=W^0+n\).

\paragraph{Inductive step.} Assume \(\tilde{W}^b \geq W^b+n\) for some \(b \geq 0\). We show \(\tilde{W}^{b+1} \geq W^{b+1}+n\) by considering two cases:

\textbf{Case 1 (Batch processed):} If \(W^b+X^b\geq n\), the queue has accumulated enough prompts to process a batch, so \(W^{b+1}=W^b+X^b-n\). Then
$$\tilde{W}^{b+1}=\max\{ 2n, \tilde{W}^b + X^b - n \}\geq \tilde{W}^b + X^b - n\geq( W^b+n )+ X^b - n=W^{b+1}+n.$$

\textbf{Case 2 (Stuck iteration):} If \(W^b+X^b<n\), insufficient prompts are available and the system cannot process a batch. The queue simply accumulates arrivals: \(W^{b+1}=W^b+X^b\). The safety buffer ensures dominance:
$$\tilde{W}^{b+1}=\max\{ 2n, \tilde{W}^b + X^b - n \}\geq 2n>W^{b+1}+n.$$
By induction, \(\tilde{W}^b \geq W^b+n\) for all \(b \in \mathbb{N}\).

\subsection{Proof of Lemma~\ref{lemma::multiple-type coupled dominace}}\label{appendix_lemma::proof of multiple-type coupled dominace}
The lemma compares the type-\(j\) entry queue in the periodic comparison system with a reflected random walk that always subtracts \(n_j\) after adding one slot of arrivals, but never falls below the safety level \(2n_j\). Fix type \(j\) and write \(\bar W^b=\bar W^b_{(j)}\), \(\widetilde W^b=\widetilde W^b_{(j)}\), \(n=n_j\), and \(\bar X^b=\bar X^b_{(j)}\). We prove by induction that \(\widetilde W^b\ge \bar W^b\).

The claim holds at \(b=0\), since \(\widetilde W^0=2n\) and \(\bar W^0=0\). Suppose \(\widetilde W^b\ge \bar W^b\). If \(\bar W^b+\bar X^b\ge n\), then the comparison queue serves \(n\) prompts and
\[
    \bar W^{b+1}=\bar W^b+\bar X^b-n.
\]
Using the induction hypothesis,
\[
    \widetilde W^{b+1}
    =\max\{2n,\widetilde W^b+\bar X^b-n\}
    \ge \widetilde W^b+\bar X^b-n
    \ge \bar W^{b+1}.
\]
If \(\bar W^b+\bar X^b<n\), then no type-\(j\) service occurs in the comparison queue and \(\bar W^{b+1}=\bar W^b+\bar X^b<n\). In this case \(\widetilde W^{b+1}\ge2n>\bar W^{b+1}\). Thus the domination holds for all \(b\), and the argument applies independently to every type \(j\).

\subsection{Nested WAIT Coupling Bound}\label{appendix_lemma::proof of nested coupled process}
This lemma is used when a downstream boundary queue receives binomially thinned arrivals from the previous segment. We index the queue after the threshold review at the start of an opportunity in which the previous segment is processed. If the downstream boundary already has a threshold cohort, Nested WAIT removes that cohort in the same prefix batch; newly surviving prompts generated by the previous segment are treated as fresh boundary input for the next review.

Let \(W^b_{(k)}\) be the residual resident boundary queue for segment \(k\ge2\) after this review, and let \(Y^b_{(k)}\) be the number of newly surviving prompts at the next boundary-crossing opportunity. The queue satisfies
\[
    W^{b+1}_{(k)}
    \le
    W^b_{(k)}+Y^b_{(k)}
    -
    n_k{\bf 1}\{W^b_{(k)}+Y^b_{(k)}\ge n_k\}.
\]
Define \(X^b_{(k)}=Y^b_{(k)}-n_k\) and
\[
    \tilde{W}_{(k)}^0=n_k,\qquad
    \tilde{W}_{(k)}^{b+1}
    =
    \max\{n_k,\tilde{W}_{(k)}^b+X^b_{(k)}\}.
\]
We prove by induction that \(\tilde{W}^b_{(k)}\ge W^b_{(k)}\). At \(b=0\), the domination holds because \(\tilde W^0_{(k)}=n_k\) and \(W^0_{(k)}=0\). Suppose it holds at \(b\). If \(W^b_{(k)}\ge n_k\), then
\[
    W^{b+1}_{(k)}
    \le W^b_{(k)}+Y^b_{(k)}-n_k
    \le \tilde{W}^b_{(k)}+X^b_{(k)}
    \le \tilde{W}^{b+1}_{(k)}.
\]
If \(W^b_{(k)}<n_k\) but \(W^b_{(k)}+Y^b_{(k)}\ge n_k\), a threshold cohort is removed at the next boundary review, and
\[
    W^{b+1}_{(k)}
    \le W^b_{(k)}+Y^b_{(k)}-n_k
    \le \tilde{W}^b_{(k)}+X^b_{(k)}
    \le \tilde{W}^{b+1}_{(k)}.
\]
If \(W^b_{(k)}+Y^b_{(k)}<n_k\), then
\[
    W^{b+1}_{(k)}=W^b_{(k)}+Y^b_{(k)}<n_k\le\tilde{W}^{b+1}_{(k)}.
\]
Thus the domination holds for all \(b\in\mathbb{N}\).

\subsection{Martingale Root Bound for Nested WAIT}\label{appendix_lemma::proof of solution to martingale exists}
This auxiliary calculation records the exponential tilting parameter used in the high-probability memory bound. Let \(0<p_k<n_k/n_{k-1}<1\), let \(\bar Y\sim\mathrm{Binomial}(n_{k-1},p_k)\), and set \(X=\bar Y-n_k\). The moment-generating equation
\[
    \mathbb{E}[e^{\theta X}]
    =
    e^{-\theta n_k}(1-p_k+p_ke^{\theta})^{n_{k-1}}
    =
    1
\]
has a unique positive solution \(\theta_k\). This is the standard Cramer--Lundberg root for a negative-drift random walk with bounded positive jumps \citep[see, e.g.,][]{asmussen2003applied}. With \(S^i_{(k)}=\sum_{r=1}^i X^r_{(k)}\), the process \(e^{\theta_k S^i_{(k)}}\) is therefore a mean-one martingale, and Doob's maximal inequality \citep[see, e.g.,][]{durrett2019probability} yields
\[
    \mathbb{P}\left(\max_{0\le i\le r}S^i_{(k)}\ge x\right)
    \le e^{-\theta_k x}.
\]
We will use the corresponding finite-horizon bound for the reflected process
\[
    Z^0_{(k)}=0,\qquad
    Z^{b+1}_{(k)}=\max\{0,Z^b_{(k)}+X^b_{(k)}\}.
\]
If \(\max_{0\le b\le B}Z^b_{(k)}\ge x\), then some excursion of the random walk has gained at least \(x\): there are \(0\le a<b\le B\) such that \(\sum_{r=a}^{b-1}X^r_{(k)}\ge x\). Applying the preceding maximal inequality from each possible excursion start and unioning over at most \(B\) starts gives
\[
    \mathbb{P}\left(\max_{0\le b\le B}Z^b_{(k)}\ge x\right)
    \le B e^{-\theta_k x},\qquad B\ge1.
\]
For completeness, we also record the elementary drift-to-root bound used in the theorem. Writing \(D=n_k-n_{k-1}p_k>0\) and \(g(\theta)=\log\mathbb{E}[e^{\theta X}]\), we have \(g(0)=0\), \(g'(0)=-D\), and \(g''(\theta)\le n_{k-1}/4\). Applying Taylor's theorem at the positive root \(g(\theta_k)=0\) gives
\begin{align}
    \theta_k \geq \frac{8 (n_k - n_{k-1} p_k)}{n_{k-1}}.\label{eq::theta lower bound}
\end{align}
